\section{Model Description}

\subsection{Neutronics models}

This work employs the OpenMC Monte Carlo particle transport software \cite{romano_openmc:_2015}
with the ENDF/B-VII.1 nuclear data library \cite{chadwick_endf/b-vii.1_2011} to generate reference
solutions and group constants for the multigroup methods. Table
\ref{table:var} details how the methods in OpenMC and Moltres treat position, direction
of travel, neutron energy, and angle-dependence in $\Sigma_s$. The hybrid $S_N$-diffusion method
corresponds to the $S_N$ method in the correction subregion and the neutron diffusion method
elsewhere.

All multigroup methods ran at a reference temperature of 900~K with the eight neutron energy group
structure developed by Jaradat for \gls{MSRE} analysis \cite{jaradat_development_2021}.
Table \ref{table:energy-group} shows the upper neutron energy bounds of the eight-group structure.
This group structure provides an appropriate balance between solution fidelity and computational
cost compared to the four- to sixteen-group structures investigated by Jaradat.

Reference \gls{MSRE} specifications \cite{robertsonMSREDESIGNOPERATIONS1965b, fratoni_molten_2020} provide the
material compositions and densities of all components in the \gls{MSRE} models.
Table \ref{table:2d-salt-composition} lists the molten salt composition for the 2-D models. The
$^{235}$U loading of 65.25 kg is the amount at which the \gls{MSRE} first achieved criticality. The
control rods consist of \gls{MSRE} 70 wt\% Gd$_2$O$_3$-30 wt\% Al$_2$O$_3$ mixture.
%
\begin{table}[!t]
  \centering
  \footnotesize
  \caption{Variable representations in OpenMC under continuous-energy (OpenMC-CE) and multigroup
  (OpenMC-MG) modes, and in Moltres under $S_N$ neutron transport and neutron diffusion. FE refers to
  finite-element discretization.}
  \begin{tabular}{l l l l l}
    \toprule
    Variable & OpenMC-CE & OpenMC-MG & $S_N$ Transport & Diffusion \\
    \midrule
    Position, $\bm{r}$ & Continuous & Continuous & FE & FE \\
    Direction of travel, $\bm{\hat{\Omega}}$ & Continuous & Continuous & Collocation & N/A \\
    Energy, $E$ & Continuous & Multigroup & Multigroup & Multigroup \\
    Angle resolution in $\Sigma_s$ & Continuous & $L=2$ & $L=2$ & N/A \\
    \bottomrule
  \end{tabular}
  \label{table:var}
\end{table}
%
\begin{table}[!t]
  \centering
  \caption{Neutron energy group structure in this work. Originally devised by Jaradat
  \cite{jaradat_development_2021}.}
  \begin{tabular}{r S}
    \toprule
    Group & {Upper energy bound [eV]} \\
    \midrule
    1 & 2.000$\times 10^7$ \\
    2 & 1.353$\times 10^6$ \\
    3 & 6.734$\times 10^4$ \\
    4 & 9.118$\times 10^3$ \\
    5 & 1.487$\times 10^2$ \\
    6 & 4.000$\times 10^0$ \\
    7 & 6.250$\times 10^{-1}$ \\
    8 & 8.000$\times 10^{-2}$ \\
    \bottomrule
  \end{tabular}
  \label{table:energy-group}
\end{table}
%
\begin{table}[!t]
  \small
  \centering
  \setlength\tabcolsep{4pt}
  \caption{\gls{MSRE} molten salt composition when the $^{235}$U loading was at 65.25 kg.}
  \begin{tabular}{l S S S S S S S}
    \toprule
    Element & {Li} & {Be} & {Zr} & {Hf} & {$^{234}$U} & {$^{235}$U} & {$^{236}$U} \\
    \midrule
    Composition [kg] & 507.27 & 293.96 & 513.97 & 0.0029 & 0.67 & 65.25 & 0.27 \\
    \bottomrule
  \end{tabular}
  \begin{tabular}{l S S S S S S}
    \toprule
    Element & {$^{238}$U} & {Fe} & {Cr} & {Ni} & {O} & {F} \\
    \midrule
    Composition [kg] & 141.91 & 0.75 & 0.13 & 0.14 & 2.27 & 3103.22 \\
    \bottomrule
  \end{tabular}
  \label{table:2d-salt-composition}
\end{table}

\subsection{1-D MSR cases} \label{sec:1d-model-setup}

We devised six 1-D slab models with increasing complexity to verify the $S_N$ and hybrid
$S_N$-diffusion implementations and to test the performance of the hybrid
method in response to various geometrical features. Figure \ref{fig:case-geom} shows the geometries of Cases 1a to
3b. All geometries have reflective boundary conditions at $x=0$~cm, thereby
forming half-core or repeating unit cell models. Cases 1a and 1b are repeating unit
cell models with reflecting boundaries on the right-side boundaries to serve as reference cases for
the base-line $S_N$ and neutron diffusion methods. Cases 2a, 2b, 3a, and 3b
are 1-D half-core models with vacuum boundaries on the right-side boundaries.
The last two cases most closely resemble the
reference \gls{MSRE} design \cite{robertsonMSREDESIGNOPERATIONS1965b}, which has centrally located control rods
and air-filled rod guide tubes.
Correction subregions for hybrid $S_N$-diffusion simulations encompass the region between $x=0$~cm
and 10~cm in order to place the right boundary more than one transport mean free path away from the
control rod and air regions. Group 1 neutrons have the longest mean free paths of 6.45~cm and
6.84~cm in the fuel and graphite regions.

\begin{figure}[htb!]
  \centering
  \includegraphics[width=0.8\columnwidth]{case-geometry}
  \caption{Geometries of the 1-D test cases. The material labeled ``mixture'' represents a
    homogeneous mixture of fuel and graphite at a ratio of 22.5\%-77.5\% by volume.}
  \label{fig:case-geom}
\end{figure}

In preliminary investigations, the original control rod material composition
(70 wt\% Gd$_2$O$_3$-30 wt\% Al$_2$O$_3$) resulted in excessively large control rod worth estimates
of approximately 55000~pcm in the 1-D models, effectively halving $k_\text{eff}$. The worth
estimates are ten times larger
than the 5500 pcm expected from full rod insertion of all three rods in the actual \gls{MSRE}
\cite{fratoni_molten_2020}. Therefore, the proportion of Gd$_2$O$_3$ was reduced to 0.35 wt\% for a
fairer assessment of rod worth in the 1-D models. This change
brought the new rod worth estimates down to approximately 17000~pcm.

\subsection{2-D MSR cases} \label{sec:2d-model-setup}

Figures \ref{fig:1/4-geom} and \ref{fig:full-geom} show the cross-sectional 2-D quarter-core and
full-core
\gls{MSRE} models and the extent of the correction subregion for the hybrid method. This work
adapted the models from the horizontal cross section of the \gls{MSRE} numerical benchmark model
\cite{fratoni_molten_2020} with the following differences in the 2-D \gls{MSRE} model for this
work:

\begin{itemize}[noitemsep]
  \item Neglected thermal expansion of all reactor components.
  \item Modeled the control rods as perfectly annular cylinders instead of beaded control elements
    strung on flexible steel cables. Neglected the inconel shells encasing the control elements and
    the steel helical cables linking the control elements together.
  \item Set channel pitch, width, and length to 2, 0.2, and 1.2 inches, respectively.
  \item Homogenized the INOR-8 (inconel alloy reactor vessel), graphite, and molten salt regions in the sample basket.
  \item Modified the outer edge of the core graphite to be circular with no jagged edges.
  \item Replaced partial channels with full-size channels near the edge of the core graphite.
  \item Ignored the thermal insulation and shield structures outside the core.
\end{itemize}

The models impose neutron vacuum boundary conditions on the outermost boundary of the INOR-8
nickel alloy vessel. The quarter-core model imposes neutron reflecting boundary conditions along the
left and bottom boundaries (Figure \ref{fig:1/4-geom}). Note that the quarter-core and
full-core models are not identical because the full-core model contains a sample basket.
For rod worth calculations, we simulated quarter- and full-core configurations with control rods
either inserted or replaced with air. Quarter- and full-core $S_N$ results were not available at
the time of writing due to an issue in the $S_N$ solver causing excessively high memory use.

For channel power distributions, we computed normalized channel fission rates as follows:
%
\begin{align}
  \text{Normalized fission rate of channel }i &=
    \frac{(\text{Fission rate in channel }i)}
    {(\text{Total fission rate in channels})/(\text{Total number of channels})} \nonumber \\
  &= \frac{\sum^G_{g=1}\int_{V_i}\Sigma_{f,g}\phi_g(\vec{r}) dV}
    {\frac{1}{N_i}\sum^{N_i}_{i=1}\sum^G_{g=1}\int_{V_i}\Sigma_{f,g}\phi_g(\vec{r}) dV}.
\end{align}
%
We subdivided the annular channel around the control rod thimble into four elements to maintain a
regular lattice distribution.
%
\begin{figure}[p]
  \centering
  \includegraphics[width=0.9\columnwidth]{quarter-core-geom-v2}
  \caption{2-D \gls{MSRE} quarter-core model and the hybrid method correction subregion.}
  \label{fig:1/4-geom}
  \vfill
  \includegraphics[width=\columnwidth]{full-core-geom-v2}
  \caption{2-D \gls{MSRE} full-core model and the hybrid method correction subregion.}
  \label{fig:full-geom}
\end{figure}

\section{Numerical Results \& Discussion} \label{sec:msre}

\subsection{1-D numerical results \& discussion} \label{sec:1d-results-sub}

\subsubsection{1-D mesh convergence tests} \label{sec:1d-mesh-conv}

We performed mesh convergence tests on Case 3b for the neutron
diffusion and $S_N$ neutron transport methods.
Mesh element sizes approximately halve with each successive refinement level up to level 4.
Figures \ref{fig:diff-mesh-k} and \ref{fig:sn-mesh-k} show the convergence of $k_\text{eff}$
estimates observed for Case 3b simulations with neutron diffusion and $S_8$ neutron transport
methods. The $k_\text{eff}$ values plateau at 0.97132 and 0.98138 for the neutron diffusion and
$S_8$ methods, confirming spatial convergence.
After two rounds of mesh refinements in Case 3b, both methods fall within
$5\times 10^{-7}$ range of the reference $k_\text{eff}$ values, below the 0.1 pcm ($10^{-5}$)
tolerance value for determining solution convergence.
%
\begin{figure}[h]
  \centering
  \begin{subfigure}[b]{0.48\columnwidth}
    \centering
    \includegraphics[width=\columnwidth]{diffusion-mesh-convergence-k}
    \caption{Neutron diffusion method}
    \label{fig:diff-mesh-k}
  \end{subfigure}
  \hfill
  \begin{subfigure}[b]{0.48\columnwidth}
    \centering
    \includegraphics[width=\columnwidth]{sn-mesh-convergence-k}
    \caption{$S_8$ neutron transport method}
    \label{fig:sn-mesh-k}
  \end{subfigure}
  \caption{Convergence of multiplication factor ($k_\text{eff}$) estimates for Case 3b across four
    levels of mesh refinement relative to the finest mesh resolution.}
  \label{fig:mesh-k}
\end{figure}

\subsubsection{$S_N$ subsolver convergence tolerance} \label{sec:conv-tol}

The hybrid method shares similarities with diffusion-based \gls{HOLO} acceleration schemes
for neutron transport methods \cite{reynolds_analysis_2023, wangDiffusionAccelerationSchemes2014}
in that power iterations are performed on the low-order model (neutron diffusion) while transport
sweeps with the latest $k_\text{eff}$ estimate are performed on the high-order model ($S_N$).
Adams and Larsen \cite{adams_fast_2002} noted that transport correction terms for \gls{HOLO}
schemes depend weakly on, and thus potentially converge faster than, the transport solution.
We investigated a range of relative convergence tolerances for the $S_N$ subsolver to 
find the optimal value for maintaining solution accuracy and minimizing computational
costs. The relative convergence tolerance for the neutron diffusion solver is fixed at $10^{-8}$.

Table \ref{table:sn-tol} lists the number of outer iterations required by the hybrid method for
a given set of $S_8$ subsolver convergence tolerance values in Case 3b.
Reducing the $S_N$ convergence tolerance generally increases outer iterations due to the increased
likelihood of updated $S_N$ boundary sources causing the initial $S_N$ residual to exceed the
convergence tolerance and triggering another flux update.

Figure \ref{fig:sn-tol} shows the $k_\text{eff}$ error estimates relative to the reference value
when $\epsilon_\text{tol}=10^{-8}$. The hybrid method exhibits superlinear convergence ($q=1.333$)
with respect to the $S_N$ subsolver convergence tolerance. These results suggest that
$\epsilon_\text{tol}=10^{-5}$ provides an optimal balance between fidelity and computational cost;
$k_\text{eff}$ error is less than 0.01 pcm and further tightening the tolerance increases the
number of outer iterations from 2 to 3.

\begin{table}[!t]
  \centering
  \small
  \caption{Number of outer iterations in hybrid method calculations of Case 3b for a given set of
  convergence tolerance values imposed on the $S_8$ subsolver.}
  \begin{tabular}{l S S S S S S}
    \toprule
    $S_8$ subsolver tolerance, $\epsilon_\text{tol}$ & {$10^{-8}$} & {$10^{-7}$} & {$10^{-6}$} & {$10^{-5}$} & {$10^{-4}$} & {$10^{-3}$} \\
    \midrule
    Number of outer iterations & 3 & 3 & 3 & 2 & 2 & 1 \\
    \bottomrule
  \end{tabular}
  \label{table:sn-tol}
\end{table}
%
\begin{figure}[!b]
  \centering
  \includegraphics[width=0.48\columnwidth]{sn-tol}
  \caption{$k_\text{eff}$ error estimates of Case 3b for a range of convergence tolerance values
  imposed on the $S_8$ subsolver relative to the reference $k_\text{eff}$ value when
  $\epsilon_\text{tol}=10^{-8}$.}
  \label{fig:sn-tol}
\end{figure}

\subsubsection{$k_\text{eff}$ and control rod worth}

Figure \ref{fig:1d-rho} shows the $k_\text{eff}$ differences of OpenMC-MG, $S_8$,
neutron diffusion, and hybrid methods relative to OpenMC-CE.
OpenMC-MG and $S_8$ methods show good agreement across all cases.
OpenMC-MG deviates from OpenMC-CE by approximately 350~pcm for Cases 2a and 3a, which indicates
that the eight neutron energy group structure introduces some discrepancy.
Increasing the Legendre polynomial order for better scattering resolution in OpenMC-MG did not
change $k_\text{eff}$ or neutron flux beyond statistical ranges. While eight neutron groups are generally
too few for high-fidelity neutronics studies, they remain within the typical two-to-ten group range
adopted in \gls{MSR} multiphysics studies \cite{krepel_dyn3d-msr_2007,
kophazi_development_2009, fiorina_modelling_2014, tiberga_results_2020}, reflecting a balance
between accuracy and computational efficiency.

\begin{figure}[htb!]
  \centering
  \includegraphics[width=\columnwidth]{rho}
  \caption{$k_\text{eff}$ differences of all neutronics methods investigated relative
  to OpenMC-CE. The standard deviation refers to the uncertainty associated with OpenMC-CE.}
  \label{fig:1d-rho}
\end{figure}

For Cases 1a and 1b, all methods show excellent agreement with OpenMC-CE as the
reactivity values fall within the one or two standard deviations. The hybrid method is consistent
with the neutron diffusion method for Case 2a, which indicates the hybrid method recovers the
diffusion solution on problems with large homogeneous regions where diffusion is largely adequate.
Both methods deviate from reference multigroup methods (OpenMC-MG and $S_8$) largely due to
neutron leakage errors on the vacuum boundary amounting to approximately -300~pcm.
Comparing Cases 2a and 3a, which differ by homogeneous and heterogeneous fuel-graphite regions, the
heterogeneity appears to induce a positive bias in $k_\text{eff}$ for the hybrid and neutron
diffusion methods. While the hybrid method continues to underestimate $k_\text{eff}$ in Cases 2b
and 3b, it shows a significant improvement over the $>1000$~pcm discrepancies logged by the neutron
diffusion method relative to the reference neutron transport methods.

For the percentage difference in rod worths in Figure \ref{fig:1d-worth}, we calculated each
method's rod worth estimates by taking reactivity
difference between the `rod out' (2a \& 2b) and `rod in' (3a \& 3b) configurations.
The neutron diffusion methods are clear outliers with rod worth errors in excess
of 8\% while the remaining multigroup methods fall within 1.8\% to 3.2\% error.
The hybrid method is consistent with OpenMC-MG and $S_8$ within 1\% error, which signals improved
control rod worth estimation over the neutron diffusion method.

\begin{figure}[h]
  \centering
  \includegraphics[width=\columnwidth]{worth}
  \caption{Percentage difference in rod worth for Cases 2 and 3 of all neutronics methods
  investigated relative to OpenMC-CE.}
  \label{fig:1d-worth}
\end{figure}

\subsubsection{Drift coefficient and flux distributions} \label{sec:flux-drift-dist}

Discussions of drift coefficient and neutron flux distributions focus on Case 3b, which resembles
the \gls{MSRE} design with a control rod inserted.
%We compare flux solutions from the deterministic multigroup methods against OpenMC-MG to
%focus on method verification outside of energy discretization errors.

\begin{figure}[p]
  \centering
  \includegraphics[width=.85\columnwidth]{case-3b-drift}
  \caption{Multigroup drift correction ($\vec{D}_g$) $x$-component distributions from the
  Moltres-hybrid and Moltres-$S_8$ solvers.}
  \label{fig:3b-drift}
\end{figure}

\begin{figure}[p]
  \centering
  \includegraphics[width=.85\columnwidth]{case-3b-flux}
  \caption{Neutron group flux distributions for Case 3b from OpenMC-CE, OpenMC-MG, Moltres-$S_8$,
  Moltres-diffusion, and Moltres-hybrid.}
  \label{fig:3b-flux}
\end{figure}

\begin{figure}[p]
  \centering
  \includegraphics[width=.85\columnwidth]{case-3b-flux-diff}
  \caption{Absolute difference in neutron group flux distributions for Case 3b from Moltres-$S_8$,
  Moltres-diffusion, and Moltres-hybrid relative to OpenMC-MG.}
  \label{fig:3b-flux-diff}
\end{figure}

Figure \ref{fig:3b-drift} shows the neutron drift distributions from reference $S_8$ and hybrid
methods for each neutron energy group for Case 3b. Reference drift data was computed
using $S_8$ flux solutions and Eq.\ \ref{eq:drift}.
We applied the adaptive buffer region boundary algorithm (described in Section
\ref{sec:hybrid-algorithm}) on the fuel region between $x=6.9375$~cm and 8.0625~cm. Due to
the regular lattice geometry and isothermal conditions, the inner boundary falls exactly at
$x=7.5$~cm, the midpoint of the fuel region.
The drift correction parameter matches the reference values within the correction subregion
from $x=0$~cm to 10~cm except near the subregion boundary at $x=10$ cm due to approximate
boundary sources derived from neutron diffusion flux solutions. This provides the justification for
imposing a buffer region within the correction subregion where drift coefficients are discarded.

Figure \ref{fig:3b-flux} shows the neutron group flux distributions from the various neutronics
methods. The discrepancies between OpenMC-CE and OpenMC-MG, particularly in groups 2, 3, 6, 7, and
8, substantiate our assessment that energy discretization errors are responsible for
multigroup $k_\text{eff}$ deviations.
While $S_8$ does not reproduce the angular resolution of groups 1 and 2 as accurately as OpenMC-MG
in the heterogeneous fuel-graphite lattice, the solutions are consistent in the lower-energy
groups. Moreover, the agreement in $k_\text{eff}$ and rod worth suggests that fine angular
resolution in the highest energy groups may be of limited importance for accurately predicting
global reactivity and heat generation in multiphysics simulations of thermal graphite-moderated
\glspl{MSR}.

Figure \ref{fig:3b-flux-diff} shows the absolute differences in neutron group flux distributions of
the Moltres-$S_8$, diffusion, and hybrid methods relative to OpenMC-MG. The neutron diffusion and
hybrid methods fare worse than $S_8$ at
resolving group 1 and 2 flux near the vacuum boundary at $x=50$~cm and group 1, 2, 5, 7, and 8
fluxes in the fuel-graphite lattice between $x=7.5$~cm and 45~cm. These flux errors contribute to
leakage and reaction rate errors responsible for their $k_\text{eff}$ discrepancies relative to
OpenMC-MG and $S_8$. The hybrid method is more
consistent with OpenMC-MG and $S_8$ between $x=0$~cm and 7.5~cm where the correction subregion
provides transport corrections, leading to better $k_\text{eff}$ and rod worth estimates over the
neutron diffusion method.

\subsubsection{Impact of correction subregion sizes}

Minimizing the correction subregion size would reduce the computational cost of the hybrid method
for time-dependent full-core simulations.
We investigated the effect of the correction subregion size on the $k_\text{eff}$ and
control rod worth estimates with Cases 3a and 3b by varying the subregion sizes from 10 cm to 40 cm
at 5 cm-intervals.

\begin{figure}[!t]
  \centering
  \begin{subfigure}[b]{0.49\columnwidth}
    \centering
    \includegraphics[width=\columnwidth]{correction-size-a-k}
    \caption{Case 3a}
    \label{fig:v1-size-a-k}
  \end{subfigure}
  \hfill
  \begin{subfigure}[b]{0.49\columnwidth}
    \centering
    \includegraphics[width=\columnwidth]{correction-size-b-k}
    \caption{Case 3b}
    \label{fig:v1-size-b-k}
  \end{subfigure}
  \caption{$k_\text{eff}$ estimates from the hybrid method for Cases 3a and 3b with different
  correction subregion sizes. Horizontal lines indicate $k_\text{eff}$ estimates from the
  OpenMC-CE, OpenMC-MG, and $S_8$ methods.}
  \label{fig:v1-size-k}
  \vfill
  \includegraphics[width=0.6\columnwidth]{correction-size-rho}
  \caption{Percentage difference in rod worth from the hybrid method relative to OpenMC-CE for
    Cases 3a and 3b with different correction subregion sizes. Horizontal lines indicate
    equivalent rod worth differences from the OpenMC-MG and $S_8$ methods.}
  \label{fig:v1-size-rho}
\end{figure}

Figures \ref{fig:v1-size-a-k} and \ref{fig:v1-size-b-k} show how $k_\text{eff}$ varies with
correction subregion size for Cases 3a and 3b. The non-monotonic $k_\text{eff}$ behavior is
consistent with
spatially non-uniform diffusion error in the heterogeneous lattice. As the subregion expands,
corrections are applied to regions with varying sensitivity to transport effects, leading to
non-monotonic error reduction.
Expanding the correction subregion does not eliminate leakage error at the outer boundary;
once the subregion spans the full domain, the hybrid method reduces to a full $S_8$ calculation
with \gls{NDA}.

Figure \ref{fig:v1-size-rho} shows the percentage difference in rod worth relative to OpenMC-CE. 
The non-monotonic dependence on correction subregion size
indicates unequal corrections to $k_\text{eff}$ provided by the hybrid method to Cases 3a and 3b.
Nevertheless, control rod worth estimates for all correction subregion sizes remain within
0.2\% of the OpenMC-MG and $S_8$ methods. These results show that the hybrid method remains largely
consistent in estimating rod worth regardless of correction subregion size. Larger correction
subregions may help with correcting transport effects arising from heterogeneity or anisotropic
scattering further away from the control rod but incurs computational expense from increased $S_N$
computation.

\subsection{2-D numerical results \& discussion} \label{sec:2d-nts-results}

\subsubsection{$S_N$ discretization convergence tests}

We investigated the impact of $S_N$ angular discretization on the neutron flux distribution and
$k_\text{eff}$ when using the hybrid $S_N$-diffusion method.

Figure \ref{fig:ray-effect} presents the flux distributions for the
first four energy groups in the central reactor region, which consists of three
air-filled thimbles and a homogenized sample basket. This region corresponds
to the correction subregion shown in Figure \ref{fig:full-geom}, excluding control rods. In group
1, the flux within the thimbles exhibits pronounced ray-effect artifacts arising from the regular
geometry and insufficient angular resolution. These artifacts diminish as the angular quadrature
order increases but remain observable at the $S_{12}$ level in groups 1 and 2. The effect becomes
progressively weaker in lower-energy groups, where shorter neutron mean free paths and increased
scattering lead to more diffusive behavior.

\begin{table}[!b]
  \centering
  \begin{minipage}{0.45\textwidth}
    \small
    \centering
    \caption{Fission cross sections in the fuel region.}
    \label{table:sigmaf}
    \begin{tabular}{l S}
      \toprule
      Group, $g$ & {$\Sigma_{f,g}$ [cm$^{-1}$]} \\
      \midrule
      1 & 0.00019 \\
      2 & 0.00011 \\
      3 & 0.00019 \\
      4 & 0.00079 \\
      5 & 0.00318 \\
      6 & 0.00316 \\
      7 & 0.01530 \\
      8 & 0.03713 \\
      \bottomrule
    \end{tabular}
  \end{minipage}
  \hfill
  \begin{minipage}{0.45\textwidth}
    \centering
    \caption{2-D full-core $k_\text{eff}$ estimates using $S_6$ to $S_{12}$
      within the hybrid $S_N$-diffusion method.}
    \label{table:2d-sn-convergence}
    \begin{tabular}{l S}
      \toprule
      Angular discretization & {$k_\text{eff}$} \\
      \midrule
      $S_6$ & 1.108414 \\
      $S_8$ & 1.108423 \\
      $S_{10}$ & 1.108422 \\
      $S_{12}$ & 1.108420 \\
      \bottomrule
    \end{tabular}
  \end{minipage}
\end{table}

From group 4 onward, the flux distributions are visually indistinguishable across all angular
discretization levels. As shown in Table \ref{table:sigmaf}, lower-energy groups exhibit larger
macroscopic fission cross sections than higher-energy groups. Because $k_\text{eff}$ is primarily
governed by thermal fission in the \gls{MSRE}, and because their flux distributions exhibit minimal
sensitivity to angular discretization,
the impact of angular resolution on $k_\text{eff}$ is negligible, as summarized in
Table \ref{table:2d-sn-convergence}.

Given the limited impact of higher angular orders on both the dominant flux shape and the
eigenvalue, we proceeded with the hybrid $S_8$-diffusion scheme for 2-D analyses as an
appropriate balance between solution fidelity and computational cost.

\begin{figure}[!t]
  \centering
  \includegraphics[width=.99\textwidth]{ray-effect}
  \caption{Group 1, 2, 3, and 4 neutron flux distributions in the
    correction subregion with $S_6$, $S_8$, $S_{10}$, \& $S_{12}$.}
  \label{fig:ray-effect}
\end{figure}

\subsubsection{2-D quarter-core results}

Table \ref{table:quarter-core} details the $k_\text{eff}$, control rod worth, and
errors relative to OpenMC-CE.
OpenMC-MG continues to overpredict $k_\text{eff}$ relative to OpenMC-CE due to energy
discretization errors. Although the hybrid method performed marginally worse than the
neutron diffusion method in the ``No Rod'' case, it significantly improves $k_\text{eff}$ estimate
in the ``Rod'' case with respect to OpenMC-MG. Consequently, the hybrid method has a 1.9\%
rod worth error relative to OpenMC-MG while the neutron diffusion method has a 17.4\% error.

\begin{table}[!t]
  \small
  \centering
  \caption{$k_\text{eff}$ and control rod worth estimates for the 2-D quarter-core \gls{MSRE}
    model. Error values are relative to OpenMC-CE.}
  \begin{tabular}{l S[table-format=1.5(2)] S S[table-format=1.5(2)] S S[table-format=4(2)] S}
    \toprule
    \multirow{2}{*}{Method} & \multicolumn{2}{c}{No Rod} & \multicolumn{2}{c}{Rod} & \multicolumn{2}{c}{Rod worth} \\
                            & {$k_\text{eff}$} & {Error [pcm]} & {$k_\text{eff}$} & {Error [pcm]} & {$\Delta\rho_\text{worth}$ [pcm]} & {Error [pcm]} \\
                            \cmidrule(r){1-1} \cmidrule(rl){2-3} \cmidrule(rl){4-5} \cmidrule(l){6-7}
	  OpenMC-CE & 1.11209(43) & {-} & 1.01740(42) & {-} & 8370(53) & {-} \\
	  OpenMC-MG & 1.11979(42) & 618 & 1.02204(41) & 446 & 8541(51) & 172 \\
      Diffusion & 1.12059 & 682 & 1.00903 & -816 & 9867 & 1484 \\
      Hybrid & 1.12174 & 773 & 1.02532 & 760 & 8383 & 13 \\
    \bottomrule
  \end{tabular}
  \label{table:quarter-core}
  \caption{Average and maximum errors in the normalized channel fission rates of
  the 2-D \gls{MSRE} quarter-core models relative to OpenMC. The mean relative standard deviation of
  OpenMC normalized channel fission rates is 0.20\%.}
  \begin{tabular}{l S S S S}
    \toprule
    \multirow{2}{*}{Method} & \multicolumn{2}{c}{No Rod} & \multicolumn{2}{c}{Rod} \\
                            & {Mean [\%]} & {Maximum [\%]} & {Mean [\%]} & {Maximum [\%]} \\
                            \cmidrule(r){1-1} \cmidrule(rl){2-3} \cmidrule(l){4-5}
    Diffusion & 0.40 & 2.63 & 2.01 & 17.44 \\
    Hybrid & 0.40 & 1.32 & 0.43 & 3.08 \\
    \bottomrule
  \end{tabular}
  \label{table:quarter-core-power}
\end{table}

Figures \ref{fig:1/4-no-rod} and \ref{fig:1/4-rod} show the channel fission rates
of the quarter-core model for the ``No Rod'' and ``Rod'' cases, respectively.
In both cases, the hybrid method reduces localized power errors near the control rod thimble
relative to the diffusion method. The introduction of the control rod exacerbates fission rate
errors in the diffusion method, with deviations reaching 2\% or more at the core center and
periphery. On the other hand, the hybrid method remained more consistent with OpenMC-CE, with
deviations largely below 1\%.

\begin{figure}[p]
  \centering
  \includegraphics[width=0.99\columnwidth]{msre-quarter-no-rod-power}
  \caption{Normalized channel fission rate distribution of the 2-D \gls{MSRE} quarter-core model
    with the rod withdrawn.
  This figure depicts the 1/8th-core distribution due to symmetry about the diagonal.}
  \label{fig:1/4-no-rod}
\end{figure}

\begin{figure}[p]
  \centering
  \includegraphics[width=0.99\columnwidth]{msre-quarter-rod-power}
  \caption{Normalized channel fission rate distribution of the 2-D \gls{MSRE} quarter-core model
    with the rod inserted.
  This figure depicts the 1/8th-core distribution due to symmetry about the diagonal.}
  \label{fig:1/4-rod}
\end{figure}

Table \ref{table:quarter-core-power} shows the average and maximum errors of the
normalized channel fission rates from the neutron diffusion and hybrid methods relative to
OpenMC-CE. Overall, the hybrid method improves channel power distribution accuracy over the
neutron diffusion method.

\subsubsection{2-D full-core results}

We perform the same $k_\text{eff}$ and fission rate distribution analyses with the 2-D full-core
model under four configurations starting from no rod inserted to all three rods inserted.
Figures of full-core fission rate distributions are
available in Appendix \ref{chap:2d-fiss-rate} and are best viewed digitally.

$k_\text{eff}$ and error estimates in Table \ref{table:full-core-k} show the same trends observed
with the 2-D quarter-core model. The neutron diffusion method performs progressively worse with
each rod added while the hybrid method remains more consistent with respect to $k_\text{eff}$
deviations across all four configurations.
Consequently, the hybrid method outperforms the neutron diffusion method by maintaining control
rod errors below 3\% relative to OpenMC-MG as shown in Table \ref{table:full-core-worth}.
Table \ref{table:full-core-power} also shows the hybrid method outperforming the neutron diffusion
method with smaller average and maximum channel fission rate errors.
Following from the symmetric quarter-core problem, the full-core results demonstrate consistent
performance by the hybrid method for asymmetric full-core geometries.

\begin{table}[htb]
  \small
  \centering
  \caption{$k_\text{eff}$ estimates for the 2-D full-core \gls{MSRE} model with the indicated rods
    inserted. Error values are relative to OpenMC-CE.}
  \setlength\tabcolsep{2pt}
  \begin{tabular}{l S[table-format=1.5(2)] S S[table-format=1.5(2)] S S[table-format=1.5(2)] S S[table-format=1.5(2)] S}
    \toprule
    \multirow{3}{*}{Method} & \multicolumn{2}{c}{No Rod} & \multicolumn{2}{c}{Rod 1} & \multicolumn{2}{c}{Rod 1 \& 2} & \multicolumn{2}{c}{Rod 1, 2 \& 3} \\
                            & {\multirow{2}{*}{$k_\text{eff}$}} & {Error} & {\multirow{2}{*}{$k_\text{eff}$}} & {Error} & {\multirow{2}{*}{$k_\text{eff}$}} & {Error} & {\multirow{2}{*}{$k_\text{eff}$}} & {Error} \\
                            & & {[pcm]} & & {[pcm]} & & {[pcm]} & & {[pcm]} \\
                            \cmidrule(r){1-1} \cmidrule(rl){2-3} \cmidrule(rl){4-5} \cmidrule(rl){6-7} \cmidrule(l){8-9}
    OpenMC-CE & 1.09859(20) & {-} & 1.06980(21) & {-} & 1.04690(18) & {-} & 1.02688(18) & {-} \\
    OpenMC-MG & 1.10637(19) & 640 & 1.07633(20) & 567 & 1.05235(20) & 495 & 1.03262(17) & 541 \\
    Diffusion & 1.10701 & 692 & 1.07120 & 122 & 1.04415 & -252 & 1.02195 & -470 \\
    Hybrid & 1.10843 & 808 & 1.07906 & 802 & 1.05553 & 781 & 1.03582 & 841 \\
    \bottomrule
  \end{tabular}
  \label{table:full-core-k}
\end{table}

\begin{table}[htb]
  \small
  \centering
  \caption{Control rod worth estimates for the 2-D full-core \gls{MSRE} with the
  indicated rods inserted. Error values are relative to OpenMC-CE.}
  \setlength\tabcolsep{5pt}
  \begin{tabular}{l S[table-format=4(2)] S S[table-format=4(2)] S S[table-format=4(2)] S}
    \toprule
    \multirow{2}{*}{Method} & \multicolumn{2}{c}{Rod 1} & \multicolumn{2}{c}{Rod 1 \& 2} & \multicolumn{2}{c}{Rod 1, 2 \& 3} \\
                            & {$\Delta\rho_\text{worth}$ [pcm]} & {Error [pcm]} & {$\Delta\rho_\text{worth}$ [pcm]} & {Error [pcm]} & {$\Delta\rho_\text{worth}$ [pcm]} & {Error [pcm]} \\
                            \cmidrule(r){1-1} \cmidrule(rl){2-3} \cmidrule(rl){4-5} \cmidrule(l){6-7}
    OpenMC-CE & 2450(25) & {-} & 4494(23) & {-} & 6357(24) & {-} \\
    OpenMC-MG & 2523(23) & 73 & 4640(24) & 146 & 6455(22) & 98 \\
    Diffusion & 3019 & 569 & 5439 & 945 & 7519 & 1162 \\
    Hybrid & 2455 & 5 & 4521 & 27 & 6323 & -34 \\
    \bottomrule
  \end{tabular}
  \label{table:full-core-worth}
\end{table}

\begin{table}[htb]
  \footnotesize
  \centering
  \caption{Average and maximum errors in the normalized channel fission rates of
  the 2-D \gls{MSRE} full-core models relative to OpenMC. The mean relative standard deviation of
  OpenMC normalized channel fission rates is 0.27\%.}
  \setlength\tabcolsep{2.5pt}
  \begin{tabular}{l S S S S S S S S}
    \toprule
    \multirow{2}{*}{Method} & \multicolumn{2}{c}{No Rod} & \multicolumn{2}{c}{Rod 1} & \multicolumn{2}{c}{Rod 1 \& 2} & \multicolumn{2}{c}{Rod 1, 2 \& 3} \\
                            & {Mean [\%]} & {Maximum [\%]} & {Mean [\%]} & {Maximum [\%]} & {Mean [\%]} & {Maximum [\%]} & {Mean [\%]} & {Maximum [\%]} \\
                            \cmidrule(r){1-1} \cmidrule(rl){2-3} \cmidrule(rl){4-5} \cmidrule(rl){6-7} \cmidrule(l){8-9}
    Diffusion & 0.45 & 2.95 & 0.94 & 12.61 & 1.35 & 15.34 & 1.67 & 17.09 \\
    Hybrid & 0.43 & 1.45 & 0.43 & 1.82 & 0.43 & 2.26 & 0.43 & 2.52 \\
    \bottomrule
  \end{tabular}
  \label{table:full-core-power}
\end{table}

Across all full-core configurations, the control rod worth and local channel power errors
consistently fall below 3\%. Global $k_\text{eff}$ estimates differ by up to 300~pcm relative to
the multigroup OpenMC-MG reference due to heterogeneous core and leakage effects.
Within the context of heterogeneous whole-core modeling for time-dependent multiphysics
simulations, these specific error margins fall within the acceptable performance bounds.
Leakage and heterogeneous effects can be mitigated in future implementations using other less
computationally intensive correction techniques. Multigroup discretization errors can also be
addressed through energy group structure optimization to
minimize the roughly 550~pcm discrepancy between OpenMC-CE and OpenMC-MG.

\subsubsection{Computational performance}

All 2-D simulations ran on two AMD EPYC Milan 7543P 32-core CPUs with 256 GB RAM each.
Across all four rod configurations, the hybrid method took an average of
451~s per 2-D full-core calculation, approximately 4.9 times longer than the neutron diffusion
method. This cost increase is expected to be representative of 3-D
calculations, as the current $S_N$ subsolver computes the angular flux for all discrete ordinates
defined on a 3-D sphere regardless of symmetries in 1-D or 2-D configurations. 
The cost remains substantially lower than that expected from full-domain $S_N$ transport for
equivalent spatial resolution. While 2-D full-core $S_8$ results are unavailable due to large
memory requirements of our current $S_N$ solver implementation, we estimate
they would take at least 100$\times$ computational expense of neutron diffusion
simulations given the $N\left(N+2\right)=80\times$ increase in problem size and strong neutron
scattering in graphite.
Data transfer times, currently taking up 25\% of the total, could be reduced through an optimal
selection of data transfer and distributed mesh partitioning options.
